\documentclass[11pt]{article}
\usepackage[T1]{fontenc}
\usepackage{lmodern,microtype,amsmath,amssymb,amsthm,mathtools}
\usepackage[margin=1in]{geometry}
\usepackage{booktabs,longtable,array,enumitem,xcolor}
\usepackage[colorlinks=true,linkcolor=blue!45!black,urlcolor=blue!45!black,citecolor=blue!45!black]{hyperref}
\hypersetup{pdftitle={Centered Pell norms and shared positivity},
 pdfauthor={Research contribution prepared for the ProveIt project},
 pdfsubject={Explicit Diophantine consistency equations and exact algebra audits},
 pdfkeywords={Diophantine equations, ZFC consistency, Pell equations, arithmetic circuits}}
\usepackage{fancyhdr}
\pagestyle{fancy}\fancyhf{}\lhead{Diophantine consistency}\rhead{Centered norms and shared positivity}\cfoot{\thepage}
\setlength{\headheight}{14pt}
\setlength{\emergencystretch}{3em}
\setcounter{tocdepth}{1}
\setlist{itemsep=3pt,topsep=5pt}
\newtheorem{theorem}{Theorem}[section]
\newtheorem{lemma}[theorem]{Lemma}
\newtheorem{proposition}[theorem]{Proposition}
\newtheorem{corollary}[theorem]{Corollary}
\theoremstyle{remark}\newtheorem{remark}[theorem]{Remark}
\newcommand{\Z}{\mathbb Z}\newcommand{\N}{\mathbb N}\newcommand{\Q}{\mathbb Q}
\newcommand{\Pos}{\operatorname{Pos}}\newcommand{\Con}{\operatorname{Con}}
\newcommand{\ZFC}{\mathrm{ZFC}}
\DeclareMathOperator{\lead}{in}
\title{Centered Pell norms and shared positivity\\\large A degree-fourteen consistency equation in twenty-two integers}
\author{Research contribution prepared for the ProveIt project\\\small AI-assisted mathematical research and executable algebra audits}
\date{23 October 2026}
\begin{document}
\maketitle
\begin{abstract}
We construct an explicit sum of ten squares in twenty-two unrestricted
integer unknowns, of exact total degree fourteen. Its integer solvability is
equivalent to finite ZFC refutability through the repository's specified
native/compiler calibration. Centering an auxiliary Pell root and applying
an elementary row operation permit the last coordinate elimination.
A second norm transformation, justified by a scale that never vanishes,
gives a degree-ten equation in twenty-four integers. The intermediate
degree-twelve equation in twenty-three integers uses only 122 binary
operations. These constructions combine the signed-power bootstrap,
a calibrated determinant, and a shared strict-positivity certificate.
Complete leading forms and exact polynomial identity checks certify the
new endpoints. Companions have thirty-one integer unknowns at degree six
and twenty-six at degree eight. Separately, adjacent-coefficient Pell descent removes an
unnecessary divisibility factor from the eleven-positive-variable form,
lowering its degree from 325 to 277 and its circuit from 113 to 111 binary
operations. A twelve-positive-variable companion has degree 78 in 103
operations. We also prove an index-elimination lemma for a different CRT
proof-certificate architecture; that lemma is not yet compiled into another
consistency polynomial. All fixed coefficient integers, polynomial circuits,
proof dependencies, checks, and source pins accompany this article.
\end{abstract}
\tableofcontents
\newpage

\section{Outcome, comparison, and the exact claim}

The smallest new low-degree construction is $P_{22}$ in
Section~\ref{sec:centered}: twenty-two unrestricted integer unknowns,
ten squared residuals, exact degree fourteen, and 123 binary arithmetic
operations. It is obtained by explicit transformations of the completely
displayed degree-twelve polynomial $P_{23}$ in Section~\ref{sec:integer}.
A degree-ten companion $P_{24}$ retains twenty-four unknowns. All
coefficient integers are materialized in the archive. Their domains are
respectively $\Z^{22}$, $\Z^{23}$ and $\Z^{24}$, with no sign conditions
left outside the equations.

The repository was inspected at commit
\texttt{2605669a532342f153b13eb6b4a0af4b54f5fa37}. Its catalogue includes
sources 31--75 and the reconciled wave-10 frontier. Comparisons in this
article refer to that pin, not to the older baseline used by each incoming
package. The source numbers below refer to directories in
\texttt{Computability/DiophantineConsistency/con-zfc-diophantine-forms}.

\begin{center}
\begin{tabular}{@{}p{3.2cm}p{5.1cm}p{6.3cm}@{}}
\toprule
Domain and objective & Prior inspected construction & This contribution\\
\midrule
Low-degree integer tradeoffs & $26$ integers, degree $16$ (68); also $24$ integers, degree $1936$ in the catalogue & $22$ integers, degree $14$; $23$ integers, degree $12$; $24$ integers, degree $10$\\
Integer sextic & $32$ integers (69, 71) & $31$ integers, eighteen residuals\\
Integer octic & $27$ integers (69, 71) & $26$ integers, fourteen residuals\\
Eleven positive integers & Degree $325$, $113$ operations (70) & Degree $277$, $111$ operations\\
Twelve positive integers & Degree $82$, $102$ operations (69) & Degree $78$, $103$ operations\\
\bottomrule
\end{tabular}
\end{center}

Each new first-row construction improves both coordinates of the prior
26-integer/degree-16 result. The 22- and 23-integer equations improve both
coordinates of the older 24-integer/degree-1936 result as well; $P_{24}$
keeps its variable count and lowers its degree. The last row trades one extra arithmetic operation
for four degrees. The nine-positive-unknown and ten-integer-variable
high-degree records are unchanged. Thus there is no assertion of an
absolute smallest equation, or of a uniform improvement in every complexity
measure. In particular, operation counts do not measure bit complexity:
some coefficient integers have more than a million bits.

\subsection{The consistency statement and its proof boundary}

Let $\operatorname{Ref}_{\ZFC}$ mean existence of a finite refutation in
the recursively specified calculus used by the pinned K215 calibration.
The target is the internal mathematical assertion
\begin{equation}
 \ZFC\vdash
 \left(\neg\exists\mathbf x\in\mathcal D\;P(\mathbf x)=0\right)
 \longleftrightarrow\Con(\ZFC).
 \label{eq:target}
\end{equation}
Here $(P,\mathcal D)$ can be any of the seven delivered endpoint/domain
pairs, in particular $(P_{22},\Z^{22})$. Consistency is syntactic:
$\Con(\ZFC)$ abbreviates absence of a
finite refutation. No assumption of consistency or arithmetic soundness
is used in the reductions.

The new coordinate transformations, positivity lemmas and Pell descent
are proved below in ordinary mathematics formalizable in ZFC. The final
interpretation uses the repository's stated native/compiler theorem
identifying its particular fixed-coefficient source with
$\operatorname{Ref}_{\ZFC}$, and its canonical-witness theorem.
Those dependencies are included and pinned, rather than replaced by finite
tests. In particular, the integer construction depends on sources 64, 71
and 75, and the positive constructions on sources 44, 61, 69 and 70.

This is an ordinary mathematical contribution with executable algebra
audits. It is not a newly compiled Lean/Rocq certificate of
\eqref{eq:target}. The repository's \texttt{OBJECT\_ZFC.md} distinguishes
external Lean equivalence from an actual derivation in its encoded ZFC
calculus; that formalization obligation is not silently declared closed
here. Likewise, \texttt{NATIVE\_CALIBRATION52.md} distinguishes the native
arithmetic from its complete formal machine/tableau bridge. The archive
retains both status documents. If one requires a new end-to-end formal
derivation file rather than the inherited ordinary calibration theorem,
that stronger deliverable remains open.

\subsection{What is new, and what is inherited}

The cone, the nearest-square positivity lemma, and the adjacent-coefficient
Pell descent already occur separately in the latest incoming packages.
We do not claim those ingredients as new. The contributions here are their
specific compositions, the additional eliminations those compositions
permit, the resulting literal smaller equations, and the proofs and exact
audits of the new endpoints. In particular, the shared guard makes
elimination of both $D$ and $V$ affordable at degree twelve. The centered
auxiliary norm, the nonzero-scale replacement of the main norm, and the
strong-norm row operation then yield the 24-integer/degree-10 and
22-integer/degree-14 endpoints. These last steps go beyond simply
substituting monically defined coordinates.
The CRT index substitution in Section~\ref{sec:homogeneous-closure} is a
separate newly developed compiler-stage improvement. No global historical
priority claim is made for any of these mechanisms.

\section{The main equation: twenty-three unrestricted integers}
\label{sec:integer}

\subsection{Fixed numerals and the complete polynomial}

Use the coherent source-34 tuple, named
\texttt{mixed-parity-unrounded-1064921}, together with source 71's two
determinant scales:
\[
 \beta=2^{1064921}-1,\qquad A_0=2^{1113},\qquad B_0=2^{666104}.
\]
The remaining three native numerals are $K_0,M_C,M_F$. Their complete
integer bytes, their hashes and their native binding are in
\texttt{integer/coefficients/BINDING.json} and its six payloads.
$M_F$ is the already shifted source mask. These symbols are fixed
coefficients, not unknowns. They are the exact integers of the pinned
source, not arbitrary integers satisfying some inequalities.

The twenty-three unknowns, all unrestricted integers, are
\begin{equation}
 J,F,\alpha,f,h,r,v,\tau,\eta,y,\gamma,o,j,E,c,a,I,u,Z,t,p_1,p_2,p_3.
 \label{eq:vars23}
\end{equation}
In the data files, $v,t,p_1,p_2,p_3$ are named
\texttt{v\_scale}, \texttt{pair\_root}, and
\texttt{pair\_pos1}, \texttt{pair\_pos2}, \texttt{pair\_pos3}.
Define only polynomial abbreviations:
\begin{align*}
 q&=\beta J+1,& C&=q-F-\alpha,&Y&=a-E,\\
 H&=4a+3,&\Delta&=a^2+H,&X&=4qu^2,\\
 D&=X+ac+\gamma H,&V&=fo-c,&k&=1-V+jc+hE,\\
 \zeta&=k-\eta,&\Pi&=1+p_1^2+p_2^2+p_3(p_3+1).
\end{align*}
Also put
\begin{align*}
 R&=\bigl((q-1)(q-F)+\alpha\bigr)(q^2-1)+(M_C+qM_F)J,\\
 G&=(A_0F)(B_0C)-(q-A_0F-B_0C)^2.
\end{align*}
The eleven residuals are
\begin{align}
 r_1&=Y-(2v+1)q^3,\nonumber\\
 r_2&=E-XY,\nonumber\\
 r_3&=c-\eta-kY,\nonumber\\
 r_4&=\tau^2+(2\tau-k)E(c-\eta)-1,\nonumber\\
 r_5&=D^2-\Delta c^2-1,\nonumber\\
 r_6&=jc-V-R,\nonumber\\
 r_7&=(K_0q+X)C-qF-Z(q-1),\nonumber\\
 r_8&=I-(2r+1)c^2,\nonumber\\
 r_9&=I^2-\Delta(f^2-1),\nonumber\\
 r_{10}&=y^2+I^2(V^2-y^2)-1,\nonumber\\
 r_{11}&=G\eta\zeta-(t^2-G-\eta\zeta+1)^2-\Pi.
 \label{eq:23residuals}
\end{align}
The proposed single Diophantine equation is
\begin{equation}
 \boxed{\displaystyle P_{23}=\sum_{\nu=1}^{11}r_\nu^2=0,
 \qquad\mathbf x\in\Z^{23}.}
 \label{eq:P23}
\end{equation}
All abbreviations are to be substituted; none introduces another
existential coordinate. Every exponent involving an unknown is fixed.
The companion JSON circuit evaluates precisely this polynomial in
122 live binary arithmetic operations.

\begin{theorem}[The 23-integer reduction]
The polynomial \eqref{eq:P23} has exact total degree twelve.
Its integer solvability is equivalent to the existence of a finite ZFC
refutation through the pinned native/compiler calibration of
sources 64, 71 and 75. Consequently its unsolvability has the consistency
interpretation \eqref{eq:target}, with the proof boundary stated in
Section 1.
\end{theorem}

We prove the new parts in detail, and specify the inherited interface at
the point where it is used.

\subsection{One certificate for two positive quantities}

For integers $b_1,b_2,b_3$ put
\[
 \Pos(b_1,b_2,b_3)=1+b_1^2+b_2^2+b_3(b_3+1).
\]
This is always strictly positive. Conversely every positive integer $n$
has this form: $4n-3$ is positive and congruent to $1$ or $5$ modulo $8$,
so Legendre's three-square theorem applies. Modulo $4$, exactly one of
the three squares is odd. Writing its root as $2b_3+1$ and the other two
as $2b_1,2b_2$ proves the assertion.

\begin{lemma}[Shared positivity, source 68]
\label{lem:shared}
For integers $x,y$, one has $x>0$ and $y>0$ if and only if there are
integers $t,b_1,b_2,b_3$ such that
\begin{equation}
 xy-(t^2-x-y+1)^2=\Pos(b_1,b_2,b_3).
 \label{eq:pair}
\end{equation}
\end{lemma}
\begin{proof}
If the left side is positive, then $xy>0$. Both numbers therefore have
the same nonzero sign. They cannot both be negative: writing
$x=-a,y=-b$, with $a,b\ge1$, makes the subtracted square at least
$(a+b+1)^2>ab$.

Conversely suppose $x,y\ge1$ and set $N=x+y-1$. Then
$xy-N=(x-1)(y-1)\ge0$. Let $m=\lfloor\sqrt N\rfloor$. Choose $t=m$
if $N-m^2\le m$, and $t=m+1$ otherwise. If $N$ is a square the error
is zero. Otherwise $|t^2-N|\le m$ and $m^2<N$, so
$(t^2-N)^2<N\le xy$. The resulting positive difference has a
$\Pos$ representation. The square root and floor select witnesses in
the proof; they are not operations in the polynomial.
\end{proof}

Apply this to $x=G$, $y=\eta\zeta$. The final row in
\eqref{eq:23residuals} therefore enforces exactly
\[
 G>0,\qquad\eta\zeta>0.
\]
It uses four private coordinates instead of the six needed by two
independent three-square representations. The equivalence holds for every
integer assignment, including assignments with a negative radix.

\subsection{The determinant and its canonical calibration}

Set $U=A_0F$ and $W=B_0C$. If $G=UW-(q-U-W)^2>0$, then $U,W$ have
one common nonzero sign and
\[
 |q-U-W|<\sqrt{UW}\le\frac{|U|+|W|}{2}.
\]
It follows that $q$ has that common sign and
\[
 |q|>\frac{A_0|F|+B_0|C|}{2}\ge |F|+|C|.
\]
Thus $F,C$ and $\alpha=q-F-C$ all have the same strict sign as $q$.
In particular $F\alpha>0$ and $FC>0$. A new zero can therefore restore
the two older native sign guards by independent $\Pos$ representations.
This implication is elementary and uses no native compiler theorem.

The converse uses canonical witnesses; the determinant is not intended
to contain every positive source assignment. Source 71 proves that the
canonical native accepting words for this exact tuple satisfy
\begin{equation}
 \frac q2\le A_0F<\frac{4q}{5},\qquad
 \frac q2\le B_0C<\frac{3q}{5},\qquad q>0.
 \label{eq:canonicalcone}
\end{equation}
Here is the numerical content of that calibration, so its dependence is
auditable. The native radix is $\rho=2^{121}$ and the cell length is
8801. Every occupied canonical cell has top $C$ digit $1$ at position
3296. Its top $F$ digit at position 8791 is
\[
 c_0=105534389824824194647382298625,
 \qquad 2^{96}\le c_0<c_0+1\le3\cdot2^{95}.
\]
The carry-free block bounds are
\begin{align*}
 \rho^{3296}&\le C_j<(1+1/\rho)\rho^{3296},\\
 c_0\rho^{8791}&\le F_j<(c_0+1)\rho^{8791}.
\end{align*}
The exponent identities
\[
 121\cdot8791+1113=1064921-97,\qquad
 121\cdot3296+666104=1064921-1
\]
give \eqref{eq:canonicalcone} after summing the canonical cells.
The full digit-support and carry proof is retained as
\texttt{integer/upstream/DETERMINANT\_GUARD.md}.

To see that these bounds imply $G>0$, write $s=q/2$,
$A_0F=s+x$, $B_0C=s+y$. Then $0\le x<3s/5$ and
$0\le y<s/5$, whence
\[
 G=s^2+s(x+y)-(x^2+y^2+xy)>\frac{12}{25}s^2>0.
\]
Source 71's bounds explicitly allow the lower dummy digits used to tune
the rank. Source 75's signed-bootstrap change modifies only its square
witness and cleared transport quotient, leaving $J,F,C,\alpha$
unchanged. Thus the canonical determinant and the signed bootstrap apply
to the same family of witnesses. This compatibility is needed for
completeness and was checked separately.

\subsection{Restoring the signed source}

Source 75 supplies a forty-integer-variable source with the arithmetic
constraints $X=4qu^2$ and
\[
 (K_0q+X)C-qF-Z(q-1)=0.
\]
Its signed Pell argument recovers $X=2^R>0$, $q>0$ and the odd rank
$R$, without assuming these signs at the start. The transport equation
then gives $q\mid Z$, restoring the older quotient. Canonical
completeness has $q=2^s$, with $s$ and $R$ odd and $R-s-2\ge0$,
so $u=2^{(R-s-2)/2}$ exists. This inherited existential source reduction
is retained in \texttt{integer/upstream/SIGNED\_BOOTSTRAP.md}.

To restore its literal arithmetic from a zero of $P_{23}$, use
\begin{align*}
 L&=q^2-1,&N&=q^3,&P&=k+E(c-\eta),\\
 C_2&=c^2,&F_2&=f^2,&g&=y^2-V^2,&B_{\rm aux}&=I^2,
\end{align*}
together with the abbreviations already displayed and the root shear
\[
 \tau_{\rm old}=\tau+E(c-\eta).
\]
All these restored coordinates are integer polynomials. The shear is
exact because, writing $P-k=E(c-\eta)$,
\[
 (\tau+P-k)^2-P^2+Pk-1
 =\tau^2+(2\tau-k)E(c-\eta)-1.
\]
Twelve of the signed source's twenty-two arithmetic rows become zero
polynomial identities; the other ten become $r_1,\ldots,r_{10}$ up to
sign. The independent audit verifies this entire identity list.

The last row restores $G>0$ and $\eta\zeta>0$ by
Lemma~\ref{lem:shared}. The determinant then restores the old
$F\alpha$ and $FC$ sign guards. This gives a signed-source zero and
therefore a finite refutation by its inherited soundness theorem.

Conversely, a finite refutation supplies canonical signed-source
witnesses. The determinant is positive on them by
\eqref{eq:canonicalcone}, and $\eta\zeta>0$. Select the four shared
positivity witnesses, apply the inverse root shear, and project away the
monically defined coordinates. Every residual in \eqref{eq:23residuals}
vanishes. This proves the two existence implications. The determinant
replacement is an existence equivalence using canonical witnesses;
the subsequent monic eliminations themselves are bijections of zero sets.

\subsection{Exact degree, not a gate-propagation estimate}

In the displayed order the residual degrees are
\[
 4,4,3,4,6,4,4,3,4,6,6.
\]
Their sum of squares has the following complete degree-twelve homogeneous
part, where every supplied unknown has degree one:
\begin{equation}
 \boxed{\lead_{12}(P_{23})=
 256\beta^4J^4u^8+I^4f^4o^4
       +\eta^4(-fo+jc+hE)^4.}
 \label{eq:leading23}
\end{equation}
Indeed the degree-six parts of $r_5,r_{10},r_{11}$ are respectively
$16\beta^2J^2u^4$, $I^2f^2o^2$, and
$-\eta^2(-fo+jc+hE)^2$. Every other residual has lower degree.
In particular
\[
 [\eta^4j^4c^4]P_{23}=1.
\]
This proves exact degree twelve after the authentic fixed integers are
substituted. It cannot be invalidated by a specialization relation among
the coefficient ports.

\subsection{The sextic and octic companions}

The elimination chain also yields useful lower-degree alternatives.
Starting from source 75's centered 34-variable sextic, replace its two
native sign guards by the determinant guard. This deletes six private
unknowns and introduces three, giving a 31-variable, 18-residual sextic.
Its coefficient $[r^2c^4]=4$ certifies the degree. Its supplied circuit
has 144 operations.

Eliminate $L=q^2-1$, $\Delta=a^2+4a+3$, and $B_{\rm aux}=I^2$.
Replace the two remaining positivity representations by the shared
certificate while retaining $\zeta$ as a coordinate. This gives
$31-3-6+4=26$ integer unknowns, fourteen residuals, degree eight and
131 operations. The private monomial $t^8$ has coefficient $1$.

Finally eliminate
\[
 \zeta=1-V+jc+hE-\eta,\quad
 D=4qu^2+ac+\gamma(4a+3),\quad V=fo-c.
\]
The first elimination makes the shared guard sextic. The other two make
the main and auxiliary norms sextic as well, so they cost no further
increase in the final degree. These are the three eliminations producing
\eqref{eq:P23}. The retained 25-variable intermediate documents this
chain, but is dominated by the 23-variable endpoint and is not counted
as a separate frontier result.

\section{Centered norms: twenty-four and twenty-two integers}
\label{sec:centered}

Two additional changes improve the variable/degree tradeoff. Both use the
integral coordinate shear
\begin{equation}
 y_{\rm old}=fo+y,\qquad y=y_{\rm old}-fo.
 \label{eq:auxshear}
\end{equation}
This is a bijection on integer coordinates, since $f,o$ are retained.
Together with $V=fo-c$, it gives the exact factorization
\begin{equation}
 y_{\rm old}^{\,2}-V^2=(c+y)(2fo+y-c).
 \label{eq:auxdiff}
\end{equation}
The leading square cancels, saving one degree in this difference.
These equations use $y$ for the centered root; it no longer denotes
the uncentered $y$ of Section~\ref{sec:integer}.

\subsection{A degree-ten equation in twenty-four integers}

Start with the 26-integer octic described above. Retain $\zeta$ as an
unknown, set $k=\eta+\zeta$, eliminate $D=X+ac+\gamma H$ and $V=fo-c$,
and apply \eqref{eq:auxshear}. The centered auxiliary residual is
\begin{equation}
 A_{\rm aux}=(fo+y)^2-I^2(c+y)(2fo+y-c)-1.
 \label{eq:aux24}
\end{equation}
Its degree is five. The directly substituted main norm
$N=(X+ac+\gamma H)^2-\Delta c^2-1$ still has degree six. A different
row representing that norm will lower its degree.

Put $b=ac+\gamma H$ and $Y=a-E$, as before, and define
\begin{equation}
 M=XE+2Eb+Y\bigl[H(2ac\gamma+H\gamma^2-c^2)-1\bigr].
 \label{eq:main24}
\end{equation}
Since $\Delta=a^2+H$, direct expansion gives
\[
 b^2-\Delta c^2=H(2ac\gamma+H\gamma^2-c^2).
\]
Consequently the following identity holds as an identity of polynomials:
\begin{equation}
 \boxed{M=YN+(X+2b)(E-XY).}
 \label{eq:nonzeroscale}
\end{equation}
The retained row $E-XY=0$ therefore makes $M=YN$. The other retained
row gives $Y=(2v+1)q^3$. This is never zero at an integer solution:
$2v+1\ne0$, and $q=\beta J+1\ne0$ because $\beta>1$ and $J$ is an
integer. Thus replacing $N=0$ by $M=0$ is reversible on the common
zeros of the other rows. This explicit nonvanishing argument excludes
the extra zero component that an arbitrary multiplication could introduce.
No positivity, native calibration or Pell theorem is used to establish it.

The unknowns of $P_{24}$ are the twenty-three names in
\eqref{eq:vars23}, together with $\zeta$; $y$ has the centered meaning.
Keep the definitions of $q,C,Y,H,\Delta,X,R,G,\Pi$ from
Section~\ref{sec:integer}, but now use $k=\eta+\zeta$, $V=fo-c$,
and $b=ac+\gamma H$. Its twelve residuals are
\begin{align}
 s_1&=Y-(2v+1)q^3, &
 s_2&=E-XY,\nonumber\\
 s_3&=c-\eta-kY, &
 s_4&=\tau^2+(2\tau-k)E(c-\eta)-1,\nonumber\\
 s_5&=M, &
 s_6&=k-R-1-hE,\nonumber\\
 s_7&=(K_0q+X)C-qF-Z(q-1), &
 s_8&=k+V-cj-hE-1,\nonumber\\
 s_9&=I-(2r+1)c^2, &
 s_{10}&=I^2-\Delta(f^2-1),\nonumber\\
 s_{11}&=A_{\rm aux}, &
 s_{12}&=G\eta\zeta-(t^2-G-\eta\zeta+1)^2-\Pi.
 \label{eq:24residuals}
\end{align}
The single equation is
\begin{equation}
 \boxed{P_{24}=\sum_{\nu=1}^{12}s_\nu^2=0,
 \qquad\mathbf x\in\Z^{24}.}
 \label{eq:P24}
\end{equation}
Every residual has degree at most five. The complete degree-ten part is
\begin{equation}
 \boxed{\lead_{10}(P_{24})=
 64a^4\gamma^2(c+2\gamma)^2(a-E)^2
 +4I^4f^2o^2(c+y)^2.}
 \label{eq:leading24}
\end{equation}
These are the leading squares of $M$ and $A_{\rm aux}$. In particular,
$[I^4f^2o^2y^2]P_{24}=4$, independently of the fixed ports.
The supplied circuit has 135 binary arithmetic operations.

\begin{theorem}[The 24-integer reduction]
The polynomial $P_{24}$ has exactly twenty-four unrestricted integer
unknowns and exact degree ten. Its integer solvability is equivalent to
that of the 26-integer octic and hence to the pinned finite-refutation
predicate. Its unsolvability satisfies \eqref{eq:target} for
$\mathcal D=\Z^{24}$.
\end{theorem}
\begin{proof}
The monic definitions of $D,V$ have polynomial inverse restorations.
The shear \eqref{eq:auxshear} is invertible over the integers.
Finally \eqref{eq:nonzeroscale}, together with $E=XY$ and the proved
$Y\ne0$, makes the main-norm replacement reversible. Each step preserves
existence of common integer zeros. A sum of integer squares vanishes
exactly when all its residuals vanish. The degree certificate
\eqref{eq:leading24} proves the remaining claim.
\end{proof}

\subsection{A degree-fourteen equation in twenty-two integers}

Start with $P_{23}$ rather than $P_{24}$, so that $\zeta$ is again the
abbreviation $1-V+jc+hE-\eta$. Apply the same shear to $y$. Define
\[
 Q=\Delta(f^2-1),\qquad S=I^2-Q,
 \qquad A_{\rm old}=y_{\rm old}^{\,2}
       -I^2(y_{\rm old}^{\,2}-V^2)-1.
\]
The residual $S$ is already present. Replace $A_{\rm old}$ by
\begin{equation}
 A_{\rm new}=(fo+y)^2-Q(c+y)(2fo+y-c)-1.
 \label{eq:aux22}
\end{equation}
The exact identity
\begin{equation}
 A_{\rm new}=A_{\rm old}+S(y_{\rm old}^{\,2}-V^2)
 \label{eq:row22}
\end{equation}
shows that this is an invertible elementary row operation in the presence
of $S=0$. It needs no division or nonvanishing assumption.

Now $I$ occurs only in its definition $I=(2r+1)c^2$ and in $S$.
Eliminate it, deleting its definition row. The strong row becomes
\begin{equation}
 S_{22}=(2r+1)^2c^4-\Delta(f^2-1).
 \label{eq:strong22}
\end{equation}
The complete single equation can therefore be written using the already
displayed $r_1,\ldots,r_7,r_{11}$ as
\begin{equation}
 \boxed{P_{22}=\sum_{\nu=1}^{7}r_\nu^2
              +S_{22}^{,2}+A_{\rm new}^{,2}+r_{11}^{,2}=0,
 \qquad\mathbf x\in\Z^{22}.}
 \label{eq:P22}
\end{equation}
The twenty-two unknowns are precisely \eqref{eq:vars23} with $I$
removed. In the unchanged residuals and abbreviations, $V=fo-c$ and
none depends on $y$, so the centered meaning of $y$ creates no further
replacement. The fixed coefficient tuple is unchanged.

The only degree-seven residual is $A_{\rm new}$. Its leading form is
$-2a^2f^3o(c+y)$; all other residuals have degree at most six. Hence
\begin{equation}
 \boxed{\lead_{14}(P_{22})=4a^4f^6o^2(c+y)^2,
 \qquad[a^4f^6o^2y^2]P_{22}=4.}
 \label{eq:leading22}
\end{equation}
The degree is exactly fourteen at the actual fixed coefficients.
The supplied circuit uses 123 binary operations.

\begin{theorem}[The 22-integer reduction]
The polynomial $P_{22}$ has twenty-two unrestricted integer unknowns,
ten squared residuals, and exact degree fourteen. Its integer solvability
is equivalent to that of $P_{23}$ and thus to finite ZFC refutability
through the pinned calibration. In particular \eqref{eq:target} holds
for $(P,\mathcal D)=(P_{22},\Z^{22})$.
\end{theorem}
\begin{proof}
The root shear is an integral bijection. In the common-zero system,
\eqref{eq:row22} has the inverse row operation of subtracting
$S(y_{\rm old}^{,2}-V^2)$. After this change the monic definition of
$I$ can be eliminated, with its displayed integer polynomial as inverse
restoration. These steps preserve the entire integer zero set up to the
specified coordinate maps. The previous soundness and completeness
implications for $P_{23}$ therefore apply. Equation~\eqref{eq:leading22}
proves exact degree fourteen.
\end{proof}

Direct substitution of $I=(2r+1)c^2$ into the old auxiliary norm would
raise the degree much more. The order matters: remove $I$ from that norm
using the strong row, cancel the leading square by centering, and only
then eliminate $I$. The 22-integer result and the degree-ten companion
are distinct useful tradeoffs; neither dominates the other in both
variable count and degree.

\section{Adjacent Pell descent and the eleven-positive-variable equation}
\label{sec:positive11}

The second contribution uses a different coherent tuple, the compact
source-44 numerals with $\beta=2^{590069}-1$. The same letters
$K_0,M_C,M_F$ now refer to that tuple. They are bound separately in
\texttt{coefficients/ACTIVE\_BINDING.json}. The four actual integer
payloads and their complete regeneration recipe are included.

\subsection{A rational coordinate is already integral}

The following descent is source 69's key lemma. Its application here
removes a factor previously inserted solely to prove integrality.

\begin{lemma}[Adjacent-coefficient integrality]
\label{lem:descent}
Let $Q>1$ be an integer. If $V\in\Z$, $y\in\Q$, and
\[
 QV^2-(Q-1)y^2=1,
\]
then $y\in\Z$.
\end{lemma}
\begin{proof}
Changing signs permits $V,y>0$; neither is zero. Write $y=b/m$ in
lowest terms with $m>0$. The equation implies $m^2\mid Q-1$, hence
$m\mid Q-1$. If $V=1$ then $y=1$ and $m=1$.

For $V>1$ define
\[
 V'=(2Q-1)V-2(Q-1)y,\qquad
 y'=(2Q-1)y-2QV.
\]
The first is an integer. The identity
$(2Q-1)^2-4Q(Q-1)=1$ preserves the equation. Its original form gives
$y>V$ and $QV>(Q-1)y$; the latter also follows by squaring positive
quantities, since the difference of their squares is $QV^2+Q-1>0$.
These inequalities give $-V<V'<V$. The preserved norm excludes $V'=0$.
Moreover $y'$ has the same reduced denominator $m$ as $y$ because
$2Q-1\equiv1\pmod m$. Replace both new coordinates by their absolute
values and use induction on the positive integer $V$. The descent reaches
$V=1$ without changing $m$, proving $m=1$.
\end{proof}

The adjacent coefficients and right side $1$ matter. For example,
$2\cdot1^2-4(1/2)^2=1$ has a nonintegral rational coordinate, so no
general statement for arbitrary positive coefficients is being used.

\subsection{The exact eleven-variable polynomial}

Supply precisely
\[
 j,F,\alpha,i,T,X,Y,n,\eta,\zeta,\gamma\in\Z_{>0}.
\]
All symbols in the following displays are polynomial abbreviations:
\begin{align*}
 J&=3j+1,&q&=\beta J+1,&C&=q-F-\alpha,\\
 E&=XY,&k&=\eta+\zeta,&c&=kY+\eta,\\
 a&=E+Y,&H&=4a+3,&\Delta&=a^2+H,\\
 \ell&=EkY,&I&=ic^2,&D_\tau&=\ell(\ell+k)+1.
\end{align*}
The removal of $q(q-1)$ from $I$ is the structural change from source 70.
Define
\begin{align*}
 R&=\bigl((q-1)(q-F)+\alpha\bigr)(q^2-1)+(M_C+qM_F)J,\\
 \widehat R&=q^3(q+F+\alpha),& M&=q^3(q-1),\\
 B_t&=(qK_0+X)C-qF,&D_f&=1+\Delta I^2,\\
 Q&=\Delta(D_f-1)=\Delta^2I^2,&\delta&=Q-1,\\
 V&=D_f(cT-R)-c,&\overline V&=D_f(cT+\widehat R),\\
 L&=3\overline V,&W&=1+(q-1)L,\\
 \overline Z&=M(\ell+k+n+L)+(q-1)Y-q^3WB_t.
\end{align*}
For the retained main carrier put
\begin{align*}
 v_0&=k-R-1,& S&=4qY,&U&=qS,\\
 D_*&=U(X+ac)+H(U\gamma+SX-v_0),\\
 A_1&=D_*^2-U^2(1+\Delta c^2).
\end{align*}
For the three-root norm set
\begin{align*}
 A&=M^2D_\tau,&B_f&=M^2D_f,&C_n&=M^2(QV^2-1),\\
 T_n&=\delta^2AB_f,&R_n&=\delta C_n,\\
 J_n&=\delta(\overline Z^2-A-B_f)-C_n,\\
 N_c&=J_n^2-4\bigl(T_n+R_n(A+B_f)\bigr),\\
 N_8&=N_c^2-64T_nR_n\overline Z^2.
\end{align*}
The polynomial is
\begin{equation}
 \boxed{P_{11}=N_8+A_1+N_8A_1=0.}
 \label{eq:P11}
\end{equation}
This literal division-free circuit has 111 binary operations.

\subsection{Radical recovery without a prime-support multiplier}

The universal identity
\begin{align*}
 F_8(z;a,b,c)
 &=\bigl((z^2-a-b-c)^2-4(ab+ac+bc)\bigr)^2-64abc z^2\\
 &=\prod_{\varepsilon\in\{\pm1\}^3}
   (z-\varepsilon_1\sqrt a-\varepsilon_2\sqrt b-\varepsilon_3\sqrt c)
\end{align*}
is a polynomial identity. For positive rational radicands all eight
displayed roots are real; it does not require independent square classes.

We recall source 70's elementary radical lemma. If three positive rational
radicands are pairwise distinct and no signed sum of their square roots
is zero, any rational signed sum of them makes each root rational.
Indeed let signed roots $x_1,x_2,x_3$ have nonzero rational sum $z$,
and write $a_i=x_i^2$. The quantities
\[
 u=\frac{z^2-a_1-a_2-a_3}{2},\qquad
 v=\frac{u^2-a_1a_2-a_1a_3-a_2a_3}{2z}
\]
are rational and equal to $\sum_{i<j}x_ix_j$ and $x_1x_2x_3$.
Then
\[
 x_1=\frac{a_1(a_1+u)}{a_1z+v},\qquad
 a_1z+v=x_1(x_1+x_2)(x_1+x_3)\ne0.
\]
The other two formulas are cyclic. Distinct radicands exclude the two
possible cancellations in each denominator.

Here the three radicands before multiplication by $M^2$ are
\[
 D_\tau,\qquad D_f,\qquad D_y=\frac{QV^2-1}{Q-1}.
\]
All of the following bounds hold for arbitrary positive supplied inputs,
before any equation or the inequality $C>0$ has been assumed. In
particular $c\ge3$, $a\ge2$, $\Delta\ge15$, $I\ge c^2$, and
\[
 1<D_\tau<D_f<Q.
\]
Put $b=cT-R$. If $b=0$, $V=-c$; otherwise
$|V|\ge D_f-c>1$. Thus $|V|>1$ in all cases and
$D_y=V^2+(V^2-1)/(Q-1)>V^2$.
If $D_y$ equalled either integer $D\in\{D_\tau,D_f\}$, then
$Q(D-V^2)=D-1$, impossible because $0<D-1<Q$.

Writing $t_\tau=\sqrt{D_\tau}$ and $f=\sqrt{D_f}$ gives
\[
 t_\tau<(a+1)c,\qquad f>(a+1)c^2\ge3c^2.
\]
When $b=0$, $\sqrt{D_y}<c+1<f-t_\tau$. When $b\ne0$,
$\sqrt{D_y}>D_f-c>2f>f+t_\tau$. One root therefore dominates
the sum of the other two, excluding a zero signed sum.

The authentic mask bounds give
$0<M_C<\beta$ and $0<M_F<2\beta$. They imply
$|R|<\widehat R$ for every positive input. For example
\[
 |R|<q^4+Fq^3+(\alpha+2)q^2<q^3(q+F+\alpha),
\]
using $q\ge16$, $\alpha\ge1$. The positive integer gap, together with
$D_f>c$, gives $|V|<\overline V$. Hence
\[
 f<\overline V,\qquad \sqrt{D_y}<2|V|<2\overline V,
 \qquad L>f+\sqrt{D_y},\quad L>k.
\]
These are exactly the estimates needed for source 70's root transport;
none uses the removed factor $q(q-1)$ in $I$.

At a zero of $N_8$, the identity
\[
 N_8=\delta^4 F_8(\overline Z;M^2D_\tau,M^2D_f,M^2D_y)
\]
supplies signed roots with
$\overline Z/M=\tau_*+f_*+y_*$. The radical lemma makes them rational.
The first two are integral because their squares are integers. The third
satisfies
\[
 QV^2-(Q-1)y_*^2=1.
\]
Lemma~\ref{lem:descent} therefore gives $y_*\in\Z$ directly.
This replaces source 70's prime-support argument and is why the stronger
divisibility of $I$ is unnecessary.

\subsection{Soundness and constructive completeness}

For soundness recover
\[
 f_{\rm old}=|f_*|,\quad u_{\rm old}=|y_*|-V,
 \quad r_{\rm old}=n+L-f_*-y_*,\quad i_{\rm old}=i.
\]
The inequalities above make $u_{\rm old},r_{\rm old}$ strictly positive.
The signed roots must be used in $r_{\rm old}$; replacing them there by
absolute values would change the transport identity. With
\[
 Z=M(\ell+k+r_{\rm old})+(q-1)Y-q^3WB_t
\]
one obtains $Z=M\tau_*$. The four source-70 weighted-carrier equations
are now
\begin{align*}
 Z^2-M^2D_\tau&=0,&A_1&=0,\\
 f_{\rm old}^2-\Delta(i_{\rm old}c^2)^2-1&=0,\\
 (V+u_{\rm old})^2-\Delta(f_{\rm old}^2-1)
            u_{\rm old}(2V+u_{\rm old})-1&=0.
\end{align*}
Its theorem applies to every positive integer transport weight, here
$L$. It restores the normalized positive source and hence a finite
refutation.

For clarity, this inherited weighted-carrier theorem has a specific
order of recovery. First $M\mid Z$ yields $q^3\mid Y$ and
$(q-1)\mid B_t$, and the positive transport gap gives $B_t>0$, thus
$C>0$. Only then are the native bounds $3q<R<q^4-2$ used. The adaptive
main carrier gives a positive main Pell root, the strong/auxiliary
equations recover its index, and its large modulus recovers $q\mid X$
and the deleted index quotient. This order avoids using native bounds
before their hypotheses have been established.

Conversely choose canonical normalized source witnesses for a finite
refutation. A positive Pell pair with ordinate divisible by $c^2$ is
available; it follows either from the canonical source or from a positive
power of the Pell matrix that is the identity modulo $c^2$. Thus
$I=ic^2$ with $i>0$. The actual auxiliary reconstruction lemma of source
61 applies to every positive such pair with $c\mid I$, canonical
$R\ge3$, $R\equiv3\pmod4$, and $c=\psi_{a+2}(R)$, supplying
positive $T,u$ and $y=V+u$. Here $\psi_A(R)$ is the coefficient of
$\sqrt{A^2-1}$ in $(A+\sqrt{A^2-1})^R$.
Choose these before computing $L,W$. Canonical transport has
$t_0=B_t/(q-1)\ge q$, $s_0=Y/q^3<k$, and $\ell<\tau<\ell+k$.
Its positive gap can therefore be set to
\[
 r=Wt_0+\tau-\ell-k-s_0>qW-2k>L.
\]
Now $n=r-L+f+y>0$ and $\overline Z=M(\tau+f+y)$, so $N_8=0$.
The adaptive main-carrier witness is supplied by the same canonical
construction. This proves completeness with no extra strong-index
congruence.

Finally $N_8$ is a square modulo $4$, and so is $A_1$ because $4\mid U$.
The equation $P_{11}=0$ is $(1+N_8)(1+A_1)=1$. The negative-unit case
would require both residuals to be $-2$, impossible modulo $4$.
Thus this single equation is equivalent to their simultaneous vanishing.

\subsection{Exact degree and its coefficient certificate}

The exact degrees are
\[
\begin{array}{c|rrrrrrrr}
 \text{register}& I&D_f&Q&V&L&\overline Z&N_8&A_1\\\hline
 \text{degree}&5&14&18&18&18&24&264&13.
\end{array}
\]
The degree-fourteen terms in $A_1$ cancel; the exact audit checks this
identity rather than trusting naive gate degrees. Writing
$q_0=3\beta j$ and $k_0=\eta+\zeta$, the leading forms satisfy
\[
 \lead_{264}(N_8)=Q_0^4\overline Z_0^8,
\]
where
\begin{align*}
 Q_0&=i^2k_0^4X^4Y^8,\\
 B_{t,0}&=(K_0q_0+X)(q_0-F-\alpha)-q_0F,\\
 \overline Z_0&=-3i^2q_0^7k_0^4X^2Y^6(q_0+F+\alpha)B_{t,0},\\
 \lead_{13}(A_1)&=32q_0^2k_0X^2Y^4
                    \bigl(4q_0^2Y\gamma+q_0^3(q_0-F)\bigr).
\end{align*}
In particular the degree-277 monomial has coefficient
\[
 [j^{60}i^{24}\eta^{49}X^{42}Y^{85}\alpha^{16}\gamma]P_{11}
 =2^7\,3^8(3\beta)^{60}>0.
\]
The degree falls by $48$, or about $14.77\%$, from source 70's 325.
The live arithmetic count falls by two multiplications, from 113 to 111.

\section{A twelve-positive-variable companion of degree 78}

This construction composes source 69's norm merger with source 70's
adaptive main carrier. Supply
\[
 j,F,\alpha,t,T,X,Y,r,\eta,\zeta,n,\gamma\in\Z_{>0}.
\]
Use $J=6j+1$, $q=\beta J+1$, and the same compact tuple and common
registers $C,E,k,c,a,H,\Delta,\ell,R,M,B_t$ as above. Set
\begin{align*}
 W&=(q-1)k+1,\\
 Z&=M(\ell+k+r)+(q-1)Y-q^3WB_t,\\
 A_0&=Z^2-M^2\bigl(\ell(\ell+k)+1\bigr).
\end{align*}
Use the adaptive $S=4qY$, $U=qS$, $D_*$ and $A_1$ from the preceding
section, in place of source 69's scales $q^4,q^5$.
The remaining registers are
\begin{align*}
 f&=t+1,&\rho&=f^2-1,&d&=\Delta c^4,&Q&=\Delta\rho,\\
 V&=f^2(cT-R)-c,\\
 h(z)&=(1-Q)z^2+2(1-Q)Vz+V^2-1,\\
 b_N&=2(1-Q)(n+V),&A_N&=d\,h(n)+(1-Q)\rho,\\
 N&=A_N^2-d\,b_N^2\rho.
\end{align*}
The single equation is
\begin{equation}
 \boxed{P_{12}=(1+A_0)(1+A_1)(1+N)-1=0.}
 \label{eq:P12}
\end{equation}

\subsection{The merger and its exact recovery}

Assume $d,\rho>0$, $Q>1$, $V>1$, $n>0$ and $N=0$.
Then $b_N\ne0$. Define the proof-only rational numbers
\[
 i_*=\frac{A_N}{d b_N},\qquad u_*=n-i_*.
\]
One gets $d i_*^2=\rho$ and $h(u_*)=0$. Writing $y_*=V+u_*$,
the latter equation is $QV^2-(Q-1)y_*^2=1$.
Lemma~\ref{lem:descent} gives $y_*\in\Z$, hence $u_*,i_*\in\Z$.
Recover the positive witnesses $i=|i_*|$ and $u=|y_*|-V$;
the second is positive because $|y_*|>|V|$.
Conversely, if $d i^2=\rho$ and $h(u)=0$, set $n=i+u$. The identity
\[
 N\big|_{\rho=di^2}=d^2h(n-i)h(n+i)
\]
shows $N=0$. This is an exact polynomial identity; it is not an
unjustified rational-root resultant inference.

The required hypotheses precede this recovery. The equations $A_0=A_1=0$
alone give $C>0$, $q^3\mid Y$ and positive first and main Pell indices.
The first index is at least two; otherwise $k=2$, $\eta=\zeta=1$
would contradict the smallest possible main ordinate. Therefore
\[
 c>kY>8XY^3\ge8q^9>2R.
\]
Since $f=t+1\ge2$, this gives
$V=f^2(cT-R)-c>c>1$. The adaptive carrier has the same hypotheses and
rank consequences here as in source 70. In particular $q\equiv2\pmod3$
and $q^3\mid Y$ make its relevant gcd at most $3$, while its remaining
modulus exceeds $R+1$. Thus replacing $q^5,q^4$ by $4q^2Y,4qY$
does not weaken the source reconstruction.

The recovered $i,u$ restore the four positive source residuals, with
transport weight $k$. Conversely, canonical witnesses have
$J\equiv1\pmod6$ because the radix exponent and repunit length are odd;
they supply $t=f-1>0$ and $n=i+u>0$. The first transport and adaptive
main-carrier gaps are positive by the canonical construction. This
proves the existence equivalence with the same native source.

The radix $q$ is even, $4\mid U$, and $b_N$ is even. Hence all three
residuals $A_0,A_1,N$ are squares modulo $4$. Their unit-product
finalizer in \eqref{eq:P12} is equivalent to their separate vanishing:
any factor $-1$ would again require a residual $-2$.

\subsection{Degree and arithmetic cost}

The exact residual degrees are $15,13,50$, so the product has degree 78.
An explicit coefficient surviving fixed-port specialization is
\[
 [j^{18}F^2\eta^{11}X^{12}Y^{23}t^8n^2\gamma\alpha]P_{12}
 =1024(6\beta)^{18}>0.
\]
The new scale computation costs one more live gate than the old scale,
giving 103 operations. This is a degree improvement from 82 to 78,
about $4.88\%$, and an explicit tradeoff rather than domination of the
old circuit in every measure.

\section{A separate compiler improvement: the modulus supplies the index}
\label{sec:homogeneous-closure}

The results above improve complete equations in the existing native-compiler
family.  This section concerns a different source architecture.  It removes
one existential coordinate from the inert-prime CRT certificate of source~75.
Its factorial, binomial and variable-exponent operations have not been
compiled and counted as a new ordinary polynomial.  Thus it is a proved
compiler lemma, not another claimed Con(ZFC) endpoint.

\subsection{The homogeneous substitution}

Fix an explicit polynomial
\(R(i,w_1,\ldots,w_m;\mathbf x)\in\mathbb Z[i,\mathbf w,\mathbf x]\),
where \(m\ge1\), and let \(R_{\rm abs}\) replace each coefficient by its
absolute value.  The external parameters \(\mathbf x\) are natural numbers.
For \(n\ge1\) and \(U\ge0\), put
\begin{align*}
 H&=R_{\rm abs}(n,U,\ldots,U;\mathbf x)+n+U+1,&
 s&=4H!,\\
 m_i&=s(i+1)-1\quad(0\le i<n),&
 B&=\prod_{i=0}^{n-1}m_i.
\end{align*}
If an arithmetic DAG already represents \(R=e-f\) with nonnegative
coefficient polynomials, the majorant \(e+f\) can be used instead, preserving
that sharing.  Write \(R=\sum_{a=0}^d R_a(\mathbf w;\mathbf x)i^a\), with
\(d=0\) for the zero polynomial, and define the ordinary integer polynomial
\begin{equation}
 \widehat R(s,\mathbf w;\mathbf x)
 =\sum_{a=0}^{d}R_a(\mathbf w;\mathbf x)(1-s)^a s^{d-a}.
 \label{eq:homogeneous-index}
\end{equation}
The right side is a finite polynomial definition.  The notation
\(s^dR((1-s)/s,\mathbf w;\mathbf x)\) explains the construction, but does
not introduce division into its polynomial representation.

For a fixed binary tree, define a norm recursively by
\(\mathcal N(z)=z\) and
\(\mathcal N(L,R)=\mathcal N(L)^2+\mathcal N(R)^2\).
Use the convention \(\binom Ak=0\) when \(0\le A<k\).

\paragraph{Homogeneous-index closure theorem.}
For each fixed \(n,U,\mathbf x\) as above, the bounded assertion
\[
 \forall i<n\;\exists\mathbf w\in\{0,\ldots,U\}^{m}
                    \quad R(i,\mathbf w;\mathbf x)=0
\]
is equivalent to the existence of natural \(A_1,\ldots,A_m\) such that
\begin{equation}
 B\mid\mathcal N\left(
   \binom{A_1}{U+1},\ldots,\binom{A_m}{U+1},
   \widehat R(s,\mathbf A;\mathbf x)\right).
 \label{eq:homogeneous-certificate}
\end{equation}
There is no separately encoded index and no index-congruence leaf.

\begin{proof}
Every prime divisor of \(m_i\) exceeds \(H\), because a prime at most
\(H\) divides \(H!\) and makes \(m_i\equiv-1\).  The \(m_i\) are pairwise
coprime: a common prime divisor of \(m_i,m_j\) would divide \(j-i\), since
it cannot divide \(s\), but \(0<|j-i|<n\le H\).  Each \(m_i\) is
\(3\pmod4\), so it has a prime divisor \(p_i\equiv3\pmod4\).
Over \(\mathbb F_{p_i}\), a sum of two squares vanishes only when both
entries vanish.  Induction on the norm tree therefore makes vanishing of
its value equivalent to vanishing of all its leaves.

Suppose \eqref{eq:homogeneous-certificate} holds.  For each \(i\), all leaves
vanish modulo \(p_i\).  The identity
\[
 (U+1)!\binom{A_j}{U+1}=\prod_{r=0}^{U}(A_j-r)
\]
and \(p_i>H>U+1\) give \(A_j\equiv w_j\pmod{p_i}\) for some
\(0\le w_j\le U\).  The same prime works for all coordinates.
Furthermore \(s(i+1)\equiv1\pmod{p_i}\), so
\[
 \widehat R(s,\mathbf A;\mathbf x)
       \equiv s^dR(i,\mathbf w;\mathbf x)\pmod{p_i}.
\]
Since \(s\) is invertible modulo \(p_i\), the residual is zero modulo
\(p_i\).  Its absolute value is at most
\(R_{\rm abs}(n,U,\ldots,U;\mathbf x)<H<p_i\), so it is zero in
\(\mathbb Z\).

Conversely, select bounded witnesses for every \(i<n\).  CRT gives
natural \(A_j\) with \(A_j\equiv w_{ij}\pmod{m_i}\).  Every \(m_i\) is
coprime to \((U+1)!\), so all binomial leaves vanish modulo \(m_i\).
The same homogeneous identity and the local equality show that the
residual leaf does also.  Pairwise coprimality then proves
\eqref{eq:homogeneous-certificate}.
\end{proof}

The proof is finite arithmetic formalizable in ZFC; it uses no consistency
or arithmetic-soundness assumption.  Existentially quantifying \(U\)
recovers unbounded witnesses at the finitely many indices.  Since the norm
has at least two leaves, it is nonnegative, and its divisibility can be
written as one equality \(\mathcal N=Bv\) with \(v\in\mathbb N\).

\subsection{The finite product and the exact saving}

As in source~75, the product defining \(B\) has the special-function
expression
\begin{align*}
 Q&=(sn)^n+1,& M&=sQ+1,\\
 B&=s^n n!\binom{Q+n}{n}\bmod M.
\end{align*}
Indeed the numerator is \(\prod_{j=1}^{n}s(Q+j)\), whose factors reduce
to \(sj-1\).  Their positive product is below \((sn)^n<Q<M\), so its
least remainder is the required product.  The variable exponent and the
special functions remain genuine compilation obligations.

Source~75 used the leaves
\[
 \binom{A_j}{U+1}\ (1\le j\le m),\qquad
 R(c,\mathbf A;\mathbf x),\qquad s(c+1)-1.
\]
Equation~\eqref{eq:homogeneous-index} removes the existential coordinate
\(c\), rather than merely dropping its check.  The number of leaves falls
from \(m+2\) to \(m+1\).  A balanced norm tree with \(k\) leaves has maximum
degree \(2^{\lceil\log_2k\rceil}\) in its leaf symbols.  When
\(m=2^r-1\), \(r\ge1\), this bound falls from \(2^{r+1}\) to \(2^r\).
This is not an unconditional degree reduction after substitution and
compilation: homogenization changes the degree in \(s\), and all graph
witnesses and shared operations still need to be counted.

More generally, a modulus of the form \(si+t\), with \(s\) a unit modulo
the selected prime divisor, supplies the residue \(i=-t/s\).  The
corresponding finite homogeneous polynomial uses \((-t)^a s^{d-a}\).
This is a reusable way of recovering an index from the modulus itself.

\subsection{Checks and limits}

The authored standard-library checker verifies 1,125 exact homogeneous
identities, 650 modular index identities, 35 product and coprimality cases,
156,584 finite-field norm evaluations, and 896 full CRT fixtures.  Of the
last set, 79 satisfy the certificate, 594 fail a range condition, and 223
fail the local residual.  Four negative controls check why one needs an
inert prime, the nested norm, a residual bound and an invertible scale.
The tests run with ordinary and optimized Python and contain no assertions
that optimization could suppress.  The universal theorem rests on the
written proof, not these finite diagnostics.

This component does not instantiate a new proof-checker source, eliminate
the special-function graphs, or establish a new minimum-variable
consistency equation independent of K215.  Its exact contribution is the
one-coordinate, one-leaf improvement to a separately identified compiler
stage.


\section{Reproducibility and interpretation of the checks}

The top-level command \texttt{python3 verify.py} checks the delivered
polynomials, their fixed numerals, and their regeneration. It requires only
Python's standard library. The PDF is compiled from this \LaTeX{} source;
the executable arithmetic does not depend on a \LaTeX{} installation.

The integer audit reconstructs the compact eleven residuals independently,
restores the arithmetic rows of the signed forty-variable source, compares
them as exact integer polynomials, and expands the final sum of squares.
It finds 16,360 monomials when the fixed ports are temporarily left as
symbolic coefficient parameters. These are not floating-point tests.
An independently written audit also verifies the exact degree certificate
after coefficient specialization: its selected coefficient is the numeral
$1$, so no relation among the fixed ports can cancel it. A separate
independent extension audit verifies the centered equations, the complete
leading forms of $P_{24}$ and $P_{22}$, and their source restorations.
It expands respectively 5,634 and 16,525 symbolic-port monomials and finds
the displayed coefficient $4$ in each leading form.

The positive audits check every named cut in the arithmetic circuits,
the universal norm identities, the cancellation in the main Pell residual,
and nonzero leading coefficients on the authentic fixed-coefficient slice.
The eleven-variable circuit has 65 multiplications, 29 additions and
17 subtractions. The twelve-variable circuit has 58 multiplications,
26 additions and 19 subtractions. Squares and scalar products are charged;
fixed leaves and reference reuse are free. Replacing a register by a
polynomial abbreviation does not create another quantified unknown.

The positive coefficient recipe is replayed from the literal transition
table and the native finite language. Missing duplicate source files in
the placed upstream package were restored byte-for-byte using their pinned
hashes, so the delivered replay works without the original checkout.
The integer component ships all six actual coefficient payloads and their
bindings, including its two determinant scales. The two components use
different coherent native tuples; their ports must never be mixed.

Finite diagnostic examples exercise the Pell descent, the shared sign
guard, and the CRT theorem, including negative controls. They support
the implementation. They cannot prove ZFC consistent, establish the absence
of all roots, or replace any universal theorem. The mathematical proof
uses the written existence implications; the checks establish exact finite
data and algebraic implementation facts.

\section{Further research questions}

\begin{enumerate}
\item \textbf{Can the 22-integer equation be lowered to degree twelve?}
The centered strong-norm row operation already eliminates
$I=(2r+1)c^2$ at degree fourteen. Its remaining bottleneck is the
degree-seven auxiliary residual, whose leading term is
$-2a^2f^3o(c+y)$. A further relation that reduces this term could improve
the endpoint without restoring a coordinate. Alternatively, can another
coordinate be removed while retaining degree fourteen?
\item \textbf{Can the determinant be recalibrated on a smaller native
tuple?} The current integer construction inherits the source-34 radix
$2^{1064921}$. The positive constructions use $2^{590069}$. Reusing smaller
ports requires a whole-tuple canonical-witness proof of the determinant
margin and the odd-rank promises; replacing one numeral in isolation is
invalid.
\item \textbf{Can the eleven-variable norm be centered at its auxiliary
root without creating a sign loop?} The bound on $f+y$ drives the degree
of the transport coordinate. A bound on $f+(y-V)$ would be much smaller
when $V>0$, but soundness must establish the needed sign before using the
transport to prove $C>0$. Removing that dependence loop would be a
substantive further improvement.
\item \textbf{Can the CRT index lemma produce a complete alternative
source?} Instantiate it with an explicit proof checker, compile factorial,
binomial, power and remainder graphs, count all witnesses and operations,
and prove the actual syntactic ZFC calibration. Only then can it be
compared with the finished native equations.
\item \textbf{Can the ordinary proof be internalized modularly?}
The polynomial coordinate maps, the strict-positivity equivalence and
the adjacent-Pell descent are small, separate targets. Their internal
ZFC derivations should precede the larger native compiler calibration.
The relevant endpoint must quantify over all internal naturals, not only
external numeral instances.
\item \textbf{Which complexity measure should govern further search?}
The variables/degree pair favors $P_{22}$ and $P_{23}$ over the older
24-variable degree-1936 tradeoff. The fewest-variable objective still favors the
high-degree nine-positive or ten-integer forms. A reproducible search
should report the domain, degree, live operations and fixed-coefficient
size together, and preserve these distinct objectives.
\end{enumerate}

\begin{thebibliography}{99}
\bibitem{repo} V. Reshetnikov and ProveIt contributors,
\emph{Con(ZFC) in Diophantine form}, repository catalogue at commit
\texttt{2605669a532342f153b13eb6b4a0af4b54f5fa37}, 23 October 2026.
\href{https://github.com/VladimirReshetnikov/ProveIt/tree/2605669a532342f153b13eb6b4a0af4b54f5fa37/Computability/DiophantineConsistency/con-zfc-diophantine-forms}{Pinned repository catalogue}.
\bibitem{s44} ProveIt source 44, \emph{Complement-packed Wang943
coefficients}, \texttt{44-b4-sparse/native/PROOF.md}, with literal
coefficient recipe and inherited compiler proofs.
\bibitem{s61} ProveIt source 61, \emph{Congruence carriers}, especially
\texttt{positive/SIGNED\_INDEX\_LEMMA.md} and the normalized positive
source interface.
\bibitem{s64} ProveIt source 64, \emph{Structural reductions},
\texttt{integer41/PROOF41.md} and \texttt{PROOF44.md}.
\bibitem{s68} ProveIt source 68, \emph{Joint positivity, free Lucas
coordinates, and commutative proof certificates},
\texttt{integer\_system/PAIR\_GUARD\_PROOF.md}.
\bibitem{s69} ProveIt source 69, \emph{Weighted cones and rational Pell
descent}, \texttt{positive/PROOF.md}.
\bibitem{s70} ProveIt source 70, \emph{Residue orientation and root
transfer}, \texttt{positive/section.tex} and \texttt{eleven/PROOF.md}.
\bibitem{s71} ProveIt source 71, \emph{Calibrated determinant guards for
Diophantine ZFC consistency}, \texttt{proofs/DETERMINANT\_GUARD.md}.
\bibitem{s75} ProveIt source 75, \emph{Signed-power bootstrap and
alternative certificate interfaces}, \texttt{guards/PROOF.md} and
\texttt{certificates/INERT\_PRIME\_CLOSURE.md}.
\bibitem{jones} J. P. Jones, Universal Diophantine equation,
\emph{Journal of Symbolic Logic} \textbf{47} (1982), 549--571.
\url{https://doi.org/10.2307/2273588}.
\end{thebibliography}
\end{document}
